\section{Related Work}
\label{sec:related_work}

\textbf{Learning-Based Modeling.}
Latent dynamics modeling approaches, such as graph neural networks (GNN)~\cite{GNN_baseline} and Bidirectional LSTM (Bi-LSTM)~\cite{bi-LSTM_baseline}, reason over vertices spatially to formulate latent dynamics of DLOs for predicting their behavior in 2-D.
Unfortunately, as we illustrate in this paper through real-world experiments, these methods are unable to accurate describe the behavior of DLOs during dynamic motions in 3D.

\textbf{Physics-Based Modeling.}
Traditionally, DLOs have been modeled using physics-based approaches.
Examples include mass-spring system models~\cite{LargeStepSimulation, diffcloth, Liu:2013:FSM},  Position-Based Dynamics (PBD) models~\cite{directionalrigidity, diminishingridgidity, xpbd},  Discrete Elastic Rod (DER) models~\cite{originalDER, originalDER2, DERapp1, DERapp2}, or Finite Element Method (FEM) models~\cite{garcia2006optimized, sin2013vega, koessler2021efficient}.
Although easy to implement, the mass-spring system imposes unnecessary stiffness to enforce inextensibility, resulting in instability during simulation.
FEM models accurately modeling DLO behavior, but are impractical for real-time prediction due to their prohibitive computational costs.
PBD and DER methods have been applied in planning and control frameworks for manipulation tasks due to their computational efficiency \cite{diminishingridgidity, xpbd, DER_manipulation1, DER_manipulation2, DER_manipulation3}.
However, PBD models are sensitive to parameter selection and assume quasi-static behavior.
As illustrated in this paper, this results in low accuracy during dynamic manipulation.
Though DER models have been qualitatively validated to describe the 3D behavior of DLOs, we illustrate that their quantitative performance while describing dynamic motion is lacking.

\textbf{Differentiable Modeling Framework.} 
Differentiable models are computational frameworks where every operation is differentiable with respect to the input parameters.
Importantly, a differentiable model enables integration into a learning framework to improve learning and can enable efficient parameter auto-tuning via gradient-based optimization during training.
There are differentiable models for deformable objects, such as DiffCloth~\cite{LargeStepSimulation, diffcloth} and XPBD~\cite{xpbd}, but these models tend to suffer from numerical instability.
This makes it challenging to apply these modeling frameworks while trying to perform parameter identification for real-world DLO manipulation.
To address these limitations, this paper proposes a differentiable DER model, called DDER, which is incorporated into a learning framework to achieve numerical stability.
This results in better predictive performance for dynamic DLO behavior when compared to state-of-the-art methods and enables auto-tuning of DLO parameters using limited real-world data.


% Crucially, as shown in \Zac{add reference}, this can significantly reduce the amount of real-world data collection required when auto-tuning the parameters, but this has not been applied to DLO modeling. 
% There are differentiable modeling frameworks for deformable objects, such as DiffCloth~\cite{LargeStepSimulation, diffcloth}, which is based on the mass-spring system model. 
% Though this modeling framework is differentiable, the mass-spring system enforces inextensibility in a stiff manner that leads to numerical instabilities during simulation.
% This makes it challenging to apply this modeling framework while trying to perform parameter identification.
% To address these limitations, this paper proposes a differentiable DER model, called DDER, which is incorporated into a learning framework that enables better predictive performance when compared to state-of-the-art methods and auto-tuning of DLO parameters using limited real-world data.

